\scanpage{57}
\mspage{13}
% archive sheet number: 56
On va peut-être se borner pour la suite à des opérations de gommage de repères, ou de surmarquage, de nature spéciale, qui se traduisent, sur \deleted{les} \inserted{chacun des} torseurs $R_{a}$ ($a\in A$) \unsure{concernés|connexes} sous $\unsure{\mathbb{Z}_{a}/(r(a))|\mathbb{Z}_{a}/r(a)}\ \deleted{\unsure{\mathbb{Z}/r(a)}}$ ($\mathbb{Z}_{a}=\mathbb{Z}\wedge_{\{\pm1\}}\omega(a)$), par une opération de contraction des groupes d'opérateurs de $\mathbb{Z}_{a}/r(a)$ à $\mathbb{Z}_{a}/r'(a)$, où $r'(a)\mid r(a)$ (opération de \uline{gommage} \uline{partiel} \inserted{\uline{uniforme}} \unsure{en} $a$) ou d'extension des groupes d'opérateurs de $\mathbb{Z}_{a}/r(a)$ à $\mathbb{Z}_{a}/r''(a)$, \unsure{où|étant} $r''(a)$ \inserted{$=d(a)r(a)$} est un multiple de $r(a)$ (opération de \uline{surmarquage uniforme} en $a$, \deleted{\unsure{)}} de multiplicateur $d(a)$). On se rend compte \inserted{ici} que le gommage \uuline{partiel} est en fait un renforcement de structure, pas un affaiblissement, et inversement pour le surmarquage uniforme --- qui, au niveau combinatoire, sera défini sans donnée de structure supplémentaire (sauf bien sûr un multiplicateur $d(a)$).
